\chapter{Editing}
\label{ch:editing}

\section{Changing a cell without retyping it}
\index{editing a cell}\index{F2 key}

Press \key{F2}, or choose \menupath{Edit > Edit cell}.

The status line changes to \msg{EDIT}, the formula bar changes colour, and the
cell's existing contents appear on it with the insertion point at the end. From
there you can move about and change what is there rather than replacing it.

What appears is the cell's \emph{contents}, not its display: a formula cell
gives you the formula, and a number formatted as \lit{\$8.50} gives you
\lit{8.5}.

\section{The edit keys}
\index{keys!editing}

\begin{cctable}{Key}{Effect}{1.55in}{3.27in}
\key{$\leftarrow$} & Insertion point left one character \\
\key{$\rightarrow$} & Insertion point right one character \\
\key{Home} & To the start of the line \\
\key{End} & To the end of the line \\
\key{Backspace} & Delete the character to the left \\
\key{Delete} & Delete the character under the insertion point \\
Any printable key & Insert that character \\
\key{Return}, \key{Tab}, \keyc{Shift}{Tab}, \key{$\uparrow$}, \key{$\downarrow$} &
Store the entry and move \\
\key{Esc} & Abandon the entry \\
\end{cctable}

Any other key is ignored, and you stay in \msg{EDIT}. Nothing you can press by
accident will drop you out of an edit half-finished.

The insertion point is shown by one highlighted character on the formula bar.
If the entry is longer than the sixty-eight characters the bar can show, the bar
scrolls sideways to keep the insertion point visible.

\section{Undo}
\index{undo}

Press \key{F5}, or choose \menupath{Edit > Undo}.

\menupath{Undo} restores the last single cell that changed --- the value it had
before your last entry or clear.

\begin{cccaution}
Undo is one cell and one level deep, and \emph{undoing is not itself
undoable}. Press \key{F5} twice and the second press does not put things back
the way they were before the first; there is nothing to restore.

Undo does not reach a pasted block, a cleared block, a sort, an inserted or
deleted row or column, an import, or a file that has been opened over the top
of your work. Sorting has its own undo --- see Chapter~\ref{ch:sorting} ---
and the others cannot be taken back at all.
\end{cccaution}

\section{Copying and pasting}
\index{copy}\index{paste}\index{clipboard}

\begin{cctable}{Key}{Command}{1.55in}{3.27in}
\keyc{Ctrl}{A} & \menupath{Edit > Select} --- start or drop a block \\
\keyc{Ctrl}{C} & \menupath{Edit > Copy} --- copies the block, or the current
cell \\
\keyc{Ctrl}{V} & \menupath{Edit > Paste} --- pastes with the cursor at the
top-left corner \\
\end{cctable}

Copying takes what the formula bar would show for each cell --- the formula, or
the unformatted number, or the text --- and pasting enters those into the sheet
as though you had typed them, with the cursor cell as the top-left corner of
what arrives. \ccprod{} says \msg{Copied} when it has taken the block.

To copy one cell, copy with no block marked. To copy a rectangle, mark it
first: see \hyperref[sec:selecting]{Selecting a block} below.

\begin{ccsteps}
\item \textbf{The format does not travel.} A cell displaying \lit{\$8400.00}
copied elsewhere arrives as \lit{8400}, in General format. The number is right;
the currency is not.
\item \textbf{References \emph{do} travel.} A pasted formula is rewritten by
the distance the block moved --- see below.
\item \textbf{A paste cannot be undone.} \key{F5} restores one cell, and a
block is not one cell.
\end{ccsteps}

The clipboard survives moving between worksheets, so a block can be copied from
one sheet and pasted into another. It does not survive quitting.

\section{Selecting a block}
\label{sec:selecting}
\index{block!selecting}\index{selection}

Press \keyc{Ctrl}{A}, or choose \menupath{Edit > Select}. That anchors one
corner of the block at the cursor. Now move the cursor: the block grows to
whatever rectangle lies between the anchor and wherever you go.

\textbf{Every movement key extends the selection} --- the arrow keys,
\key{PgUp} and \key{PgDn}, \key{Home}, \key{End}, and a click of the mouse.
There is nothing new to learn, because the keys are the ones from
Chapter~\ref{ch:moving}.

The anchor may end up below and to the right of the cursor; it makes no
difference, and the block is the rectangle between them either way.

\begin{cctable}{To}{Press}{2.4in}{2.42in}
Start a block & \keyc{Ctrl}{A} \\
Grow or shrink it & Any movement key \\
Drop it & \keyc{Ctrl}{A} again, or \key{Esc} \\
Copy it & \keyc{Ctrl}{C} \\
Clear it & \key{Del} \\
\end{cctable}

A command that acts on ``the selection'' acts on the current cell when there is
no block, so \keyc{Ctrl}{C} and \key{Del} need no block to be useful.

\section{What happens to formulas}
\index{formulas!adjusted on paste}\index{references!relative}

A pasted formula does not arrive unchanged. Its references move by the same
distance the block did, in both directions at once.

Copy \fml{=A1+B1} five rows down and two columns across, and what lands is
\fml{=C6+D6} --- which is still the two cells above it, exactly as it was
before. Column letters carry properly: one column right of \cellref{Z} is
\cellref{AA}.

To stop a reference moving, write a \lit{\$} in front of the part that should
stay:

\begin{cctable}{Written}{When the formula is copied}{1.4in}{3.42in}
\fml{A1} & Both the column and the row move \\
\fml{\$A1} & The column stays at \cellref{A}; the row moves \\
\fml{A\$1} & The row stays at 1; the column moves \\
\fml{\$A\$1} & Neither moves --- always \cellref{A1} \\
\end{cctable}

That is the ordinary spreadsheet rule, and it is what makes a column of totals
worth copying rather than retyping: put \fml{=B2*\$C\$1} in one cell, copy it
down the column, and every row multiplies its own \cellref{B} by the one rate
in \cellref{C1}.

\begin{ccnote}
The same rewriting runs when a formula is carried along by anything else that
moves it: inserting a row, deleting a row, sorting, and undoing a sort.
Chapter~\ref{ch:rows} and Chapter~\ref{ch:sorting} describe those.

Two things are deliberately left alone. A reference that names another sheet
--- \fml{Q1!B5} --- is never rewritten, and neither is any formula longer than
128 characters, which is skipped whole rather than half-adjusted.
\end{ccnote}

\section{Clearing}

\key{Backspace} or \key{Delete} in \msg{READY}, or \menupath{Edit > Clear},
empties the current cell --- or the whole block, if one is marked.

A cleared cell is genuinely empty, not zero. In a range, \fn{SUM} and
\fn{AVERAGE} skip it rather than counting it as nothing. Read on its own by a
formula, though, an empty cell reads as zero --- \fml{=Z90+1} is 1.
Chapter~\ref{ch:formularef} sets out the distinction.
